\section{模型评价}
\label{sec:evaluation}

\subsection{模型的优点与创新点}
\subsubsection{方法层面}
\textbf{把不可检验的前提显式化为条件，而不是把它留在拟合过程里。}
广义标度律的质量指数只有在数据量指数给定的前提下才可辨识，
而跨来源共享数据量指数这一前提在本数据上检验不通过。
本文没有回避这一结果，而是在给出质量指数的同时给出该前提不成立时的另一条读数，
并说明两条读数各自支持的结论，使读者能够判断结论的适用边界。

\textbf{用定义的分歧本身作为结论。}
「预算跨越量级时最优策略是否质变」在缺乏公认定义时不可判定。
本文并报三种自然定义并逐一识别，发现其中一种给出明确位置且不依赖任何人为阈值或网格步长，
另一种的位置依赖网格步长，第三种则因判据不是网格不变量而无法给出结论。
由此得到的判断是：转移存在，但形态是连续饱和而非突变。
这一结论比单一定义下的一个位置坐标更为稳健，也解释了几种定义为何不矛盾。

\textbf{以解析结构替代纯数值求解。}
优化问题的最优解在多数预算档位触及质量上界，内点条件并不刻画最优解。
本文按活跃约束集分支枚举，把内点段与边界降维段分别解析求解，
再由一阶条件的约减形式给出两个可供交叉验证的关系式，
并用中心差分重建、邻域采样与全局数值优化三条独立路径验证分支枚举的完备性。

\textbf{把不确定性拆到它真正的来源上。}
归因分解的点值在端点定义变动下会发生符号翻转，而重抽样区间在结构上不可能覆盖这一来源。
本文把端点定义与函数规格两种来源分别扰动并一并报告，
由此得出「任何具体百分比分割都不可主张」的判断，并改以方向与量级作为结论。

\subsubsection{建模层面}
\textbf{门控式质量合成。}质量指标中既有越高越好的价值类，也有越低越好的污染类。
本文用价值维度均值乘以污染维度补数的门控形式合成质量评分，
使污染指标以乘性方式施加惩罚而非在加权平均中相互抵消，
并由此自然给出「同一文本在两类维度上同时取高分」这一冲突的定义。

\textbf{逐 token 归一化处理长度混淆。}按序列求和的质量信号与文档长度高度相关，
其表观方向由长度效应主导。本文对这类字段按长度归一化后再参与合成，
使其与独立质量指标的秩相关方向由负转正，避免了因量纲错误而引入的虚假矛盾。

\textbf{依据约束结构选择求解方式。}优化问题的最优解是否落在内点由预算水平决定，
本文先以探针确定各档位的分支归属，再据此设计分支枚举，而不是默认内点方法普遍适用。

\textbf{把预测区间与纯趋势基线同表对照。}本文的预测模型在点值与区间宽度上均不优于
不含分解结构的纯趋势基线。这一比较如实列出，用以说明结构化分解的价值在于
提供可归因、可情景化的区间，而非提高点预测精度。

\subsection{模型的缺点与适用范围}
\subsubsection{数据条件带来的限制}
\begin{enumerate}[label=(\arabic*),leftmargin=2.2em,itemsep=2pt]
  \item \textbf{单一标量质量评分无法同时做到跨域可比与域去混淆。}
    本文使用的 22 个指标均为面向现代网页文本训练的判别器输出，
    其中多数指标对本批数据的领域标签具有很高的可分性。
    这是一个可证的限制：在这批指标上，任何单一标量评分都会部分地编码领域身份。
    因此本文对领域间的质量比较只在最高一档的集合层面作出，不对档内位次作断言。
  \item \textbf{配比效应的形状不可跨源验证。}附件中没有文件同时给出配比与规模、
    数据量的联合观测，配比系数只能由问题一的模型引入并作为承重前提声明；
    相应地，本文只报告各领域边际效应的方向与相对排序，不报告绝对量级。
  \item \textbf{质量指数的解释绑定在特定刻度上。}问题一给出的经验质量评分落在
    一个较窄的取值区间内，与半合成网格的设计刻度并非同一对象。
    本文不做静默的线性映射，质量指数的解释因此限定在设计刻度上。
  \item \textbf{部分结论建立在半合成数据上。}广义标度律相对经典形式的改进
    只能配备质量变化的半合成网格上验证；该结论带补充分析的性质，
    不能替代真实实验观测。
  \item \textbf{高预算档位是外推。}最高档预算对应的最优参数量与数据量
    分别达到真实观测上界的数倍与十倍以上，其结果不应被理解为可靠预测。
  \item \textbf{时间窗口不足。}问题四的密集数据仅覆盖 9.1 个月，
    而预测视界为 12 至 24 个月。预测区间偏宽是数据条件的结果，而非算法缺陷。
  \item \textbf{技术通道不可直接观测。}该通道由规模项之外的月度共同变动构成，
    同时吸收未观测的数据质量变化与模型对基准的过拟合，
    因此只能表述为非规模技术与数据工程的综合进步，不能解释为纯粹的算法改进。
  \item \textbf{预测未优于基线。}在预声明的预测判据上，纯趋势基线同样满足，
    且主方法的点值低于基线、区间宽度约为基线的数倍。
\end{enumerate}

\subsubsection{方法选择的代价}
\begin{enumerate}[label=(\arabic*),leftmargin=2.2em,itemsep=2pt]
  \item \textbf{质量与数据量的替代关系只在比例意义下成立。}
    所采用的有效数据量形式使质量与数据量的弹性之比恒等于质量指数，
    因此本文不主张二者的替代率随配置位置变化；随位置变化的只有质量与参数量之间的替代关系。
  \item \textbf{领域配比未参与联合优化。}配比项对配比是线性的，最优解必然落在单纯形顶点。
    本文因此固定配比、只报告其边际方向与排序，代价是放弃了「配比随预算变化」这一维度。
  \item \textbf{基线质量水平未经外部确认。}本文在其合理范围内做扫动并报告结论稳定，
    但该取值本身仍是一个需要外部依据的参数。
  \item \textbf{部分判据在执行后才被发现设计不足。}问题一的目标口径替换检查中，
    预声明的判据对每个领域系数恒为正，因而无法被证伪；真正有信息量的是系数向量的排序反转。
    本文如实记录这一判据设计失效，不事后更换判据以改善结论。
\end{enumerate}

\subsection{模型的推广}
上述方法中不依赖于具体题设的部分可以迁移到相邻问题。

\textbf{正负向指标并存时的合成。}门控形式适用于任何由「越高越好」与「越低越好」两类
指标构成的多维评价体系，只要两类指标在语义上可分；
此时冲突样本的识别退化为在两类维度上同时取高分的样本筛选，其阈值由分位数锚定，
不随样本量漂移。这一构造可直接用于其他语料、其他模态的数据质量治理场景。

\textbf{转移现象的识别框架。}本文对结构性转移给出的三种定义——
活跃约束集变化、成本分项占比排序变化、分配份额突变——构成一个可复用的识别框架：
当一个问题出现「最优策略是否随参数发生质变」的提问时，可按此框架分别识别，
并优先采用不依赖人为阈值与网格步长的那一种作为主定义。
该框架不限于算力配置问题，同样适用于其他资源分配与最优控制场景。

\textbf{不确定性来源的分解。}把点值的不稳定性归因到重抽样、函数规格与样本端点
三个来源，并分别扰动，是一种可以普遍采用的稳健性做法：
它避免了用重抽样区间去覆盖它结构上无法覆盖的变异来源，
也避免了在发现脆弱性之后事后放宽判据。

\textbf{跨尺度映射的误差传播。}当两个评价尺度之间只有少量高可比样本时，
把映射误差作为独立随机项并入蒙特卡洛预测，可以给出诚实的区间宽度；
这一做法可推广到任何需要在中间量尺度与目标量尺度之间转换的外推问题。

\textbf{进一步工作。}若要收紧本文的主要限制，需要补充两类数据：
同时给出配比与训练规模配置的实验，用以检验配比效应的可迁移性；
以及跨越更长窗口的模型评测记录，用以检验得分上界的饱和形态并压缩预测区间。
